\documentclass[a4paper,10pt]{report}
\usepackage[utf8x]{inputenc}

% Title Page
\title{Ganeti}
\author{aubo}


\begin{document}
\maketitle
\chapter{Ganeti project}

\section{Introduction}
\includegraphics[width=10cm,height=7cm]{images/logo_ganeti.png}

Ganeti est un un manager de cluster de machine virtuelles. Il combine la virtualization et la réplication en temps réel de disque.
Ganeti offre un plateforme de haute disponibilité

Ganeti utilse Python, Xen, KVM, DRDB, LVM et OpenSSL.
Développé par Google depuis Aout 2007

\includegraphics[width=10cm,height=7cm]{images/image1.png}

Ce Projet est sous licence GNU GPLv2.

Site du projet :

http://code.google.com/p/ganeti/


\section{Architecture}

\end{document}  



Abstraction of computer resources
 CPUs, memory, storage, network
• Advantages
 Consolidation, increase hardware utilization
 Transparent for user
 Flexibility
• Disadvantages
 Depending on application: performance losses
• Different types
 Paravirtualization
 Full virtualization
• Hypervisor
Copyright by Google Inc
3
What is Ganeti and why should you use it?
• Software to manage clusters of virtual servers
 Automation allows you to scale easily
 Makes it simple to manage 10s of nodes and 100s of instances
• Combines virtualization and data replication
 All integrated in a unified interface
 Virtual systems are portable between nodes
• Hypervisor backends
 Abstraction layer
 Currently based on Xen, but others are possible
Copyright by Google Inc
4
Terms
• Node
 Physical machine
 Xen Dom0
• Instance
 Virtual machine
 Xen DomU
• DRBD
 Distributed Replicated Block Device, http://www.drbd.org/
 Used for data replication
• LVM (Logical Volume Manager)
 Used to manage instances' volumes

